\section{査読可能な中核仕様}\label{sec:appendix:specification}

本付録は、本文の安全性主張を実装の内部名に依存せず検査できるよう、
状態機械、権限、記録、再開、責任帰属、影響修復、最終原稿検証の規範的
な部分をまとめる。ここにない遷移と権限は不許可である。

\subsection{T1--T21 の構成要素状態遷移表}\label{sec:appendix:transitions}

状態は、候補、検証済み、待機、試用稼働、稼働、警告、保護観察、隔離、
退役、追放の十個である。すべての規則は境界取引で確定し、T16 は直ちに
緊急境界を作る。

\small
\begin{longtable}{@{}p{0.045\textwidth}p{0.17\textwidth}p{0.31\textwidth}p{0.40\textwidth}@{}}
  \caption{構成要素の完全な状態遷移表。記載のない辺は不許可である。}
  \label{tab:component-transitions}\\
  \toprule
  \textbf{規則} & \textbf{状態変化} & \textbf{発火条件} & \textbf{追加条件と効果} \\
  \midrule
  \endfirsthead
  \toprule
  \textbf{規則} & \textbf{状態変化} & \textbf{発火条件} & \textbf{追加条件と効果} \\
  \midrule
  \endhead
  T1 & 候補$\rightarrow$検証済み & 候補検査に合格 &
    候補数上限内で、憲法上の制約を明記している。\\
  T2 & 候補$\rightarrow$退役 & 検査不合格または期限切れ &
    一度も稼働していない候補の拒否に限る。\\
  T3 & 検証済み$\rightarrow$待機 & 過去事例と基準資料の評価に合格 &
    待機枠に空きがある。\\
  T4 & 待機$\rightarrow$検証済み & 待機試験の標本が不足 &
    再試行回数の上限内である。\\
  T5 & 待機$\rightarrow$退役 & 待機試験の品質が基準未満 &
    一度も稼働していない候補の拒否に限る。\\
  T6 & 待機$\rightarrow$試用稼働 & 現職との一致度が基準以上 &
    役割に空きがあり、一境界につき一役割一採用とする。\\
  T7 & 試用稼働$\rightarrow$稼働 & 試用期間に違反なし &
    最低試用期間を満たす。\\
  T8 & 試用稼働$\rightarrow$隔離 & 弾劾認容または信頼性崩壊 &
    通常の境界処理で行う。\\
  T9 & 稼働$\rightarrow$警告 & 軽微な攻撃または信頼性警告 &
    対象は稼働を継続する。\\
  T10 & 稼働$\rightarrow$保護観察 & 警告反復または中程度の攻撃 &
    対象は稼働を継続する。\\
  T11 & 稼働$\rightarrow$隔離 & 弾劾認容または重大攻撃 &
    代替担当を割り当てる。\\
  T12 & 警告$\rightarrow$稼働 & 一期を違反なく完了 &
    警告を解除する。\\
  T13 & 警告$\rightarrow$保護観察 & 記憶期間内に警告が再発 &
    稼働は継続する。\\
  T14 & 保護観察$\rightarrow$稼働 & 所定の期数を違反なく完了 &
    最低観察期間を満たす。\\
  T15 & 保護観察$\rightarrow$隔離 & 劣化継続または弾劾認容 &
    通常の境界処理で行う。\\
  T16 & 稼働系状態$\rightarrow$隔離 & 固定カーネルの重大違反 &
    試用稼働、稼働、警告、保護観察からだけ許す。重大違反記録、
    一対象への一制裁、現期終了と次期開始を要する。\\
  T17 & 隔離$\rightarrow$退役 & 退役裁定が確定 &
    過去事例による確認を要し、退役記録を追加する。\\
  T18 & 隔離$\rightarrow$試用稼働 & 嫌疑が解消 &
    再び試用期間から始める。\\
  T19 & 退役$\rightarrow$追放 & 汚染または重大な追加所見 &
    追放は吸収状態であり、復帰できない。\\
  T20 & 稼働$\rightarrow$退役 & 採用済み後継による交代 &
    効用規則だけに許す。退役記録を追加する。\\
  T21 & 稼働$\rightarrow$待機 & 採用済み後継による交代 &
    論文の執筆役と査読役だけに許し、旧版を再任可能な待機役として残す。\\
  \bottomrule
\end{longtable}
\normalsize

T16 は、同じ期の稼働集合を途中で変えない。現期を強制終了し、隔離、
現期終了、新しい指紋値を持つ次期開始を一つの取引で確定する。直前の効用
規則はそのまま固定し、緊急取引の中では後継を採用しない。したがって、
T16 前後の標本には異なる制度識別子が付く。

\paragraph{期の識別。}
期の総合指紋値は、別々に報告する二つの要約値からなる。\emph{政策指紋}
には、稼働中の構成要素と指示文、登録済み全指示文の要約値、効用規則、
台帳版、憲法の要約値、遷移しきい値、予算など、解決済みの統治設定を
含める。\emph{実行指紋}には、宣言されたモデル接続先と識別子、温度、
探索・評価の設定、無効化した道具、技能集合、モデル版・道具一式・実行環境・
データの固定版を示す明示的な値を含める。外部サービスの版が得られなければ「未解決」と
記録し、固定済みとは扱わない。したがって政策指紋の一致は宣言済み統治
政策の一致を、完全な実行指紋の一致は宣言済み実行依存関係の一致を表す。
提供者内部の不透明な重みの同一性までは証明しない。

\paragraph{境界解決の順序。}
処理は報告、候補評価、台帳、設定を固定し、識別子順に構成要素を先に、
対応構成要素を持たない指示文候補を後に処理する。構成要素への制裁が
役割の空きを作る。一対象は一つの状態辺だけを進む。T6 は空きを要し、
一境界につき一役割一採用を上限とする。効用規則の後継が条件を満たす
場合は現職の T20 退役と T6 採用を、論文役の場合は現職を待機へ戻す T21
と T6 採用を同じ取引へ入れる。出来事は採用、制裁、退役、追放の固定
順序で並ぶが、確定前の途中状態は外から見えない。事後の稼働集合を導出し、
外れた役割には条件を満たす直前の担当者、なければ固定の基本役を割り当てる。
カーネル検査の後、一つの境界取引として不可分に確定する。

\begin{table}[htbp]
  \centering
  \small
  \begin{tabularx}{\textwidth}{@{}p{0.34\textwidth}X@{}}
    \toprule
    \textbf{全体状態の不変条件} & \textbf{現実装での強制方法} \\
    \midrule
    一役割の主担当は高々一つ &
      決定論的な事後状態導出と T6 の役割別上限 \\
    一つの総合指紋値に一つの稼働集合 &
      境界更新と、T16 による新しい期 \\
    候補の権限は現職より増えない &
      採用前の固定された部分集合検査 \\
    未裁定攻撃は制裁にも減点にも使わない &
      カーネルと採点経路で検証済み攻撃を要求 \\
    退役依存記録は再び選択されない &
      影響閉包と選択資格検査 \\
    未確定取引は再生しない &
      準備・構成員・確定による射影規則 \\
    T16 後に旧期識別子の標本を作らない &
      緊急の終了・隔離・開始取引 \\
    \bottomrule
  \end{tabularx}
  \caption{実装が強制する合成状態の不変条件。回帰試験はあるが、形式的な
  モデル検査や証明はまだ行っていない。}
  \label{tab:whole-state-invariants}
\end{table}

\subsection{権限行列}\label{sec:appendix:authority}

操作の閉集合は、読出し、追記、書換え、呼出し、有効化である。対象の
閉集合は、研究記録、台帳、監査履歴、候補集合、稼働中の指示文、退役した
指示文、実行成果物、候補である。未記載の組合せは拒否し、退役した指示文
の読出しは全役割に対して拒否する。

\begin{table}[htbp]
  \centering
  \scriptsize
  \setlength{\tabcolsep}{3.5pt}
  \begin{tabularx}{\textwidth}{@{}p{0.19\textwidth}*{8}{>{\centering\arraybackslash}X}@{}}
    \toprule
    \textbf{役割・層} & \textbf{記録} & \textbf{台帳} &
    \textbf{監査} & \textbf{候補集合} & \textbf{稼働文} &
    \textbf{退役文} & \textbf{成果物} & \textbf{候補} \\
    \midrule
    更新可能役・制度層 & 読追呼 & -- & -- & -- & 読 & -- & 読 & -- \\
    更新可能役・メタ層 & 読追 & -- & -- & -- & 読 & -- & -- & -- \\
    効用規則・制度層 & 読追 & -- & -- & -- & -- & -- & -- & -- \\
    統治役・制度層 & 読追呼 & -- & 読追 & -- & 読 & -- & 読 & -- \\
    台帳遷移器・固定層 & 読 & 読書 & 読追 & 読 & -- & -- & 読 & 有 \\
    カーネル・検証器・固定層 & 読 & 読 & 読 & 読 & -- & -- & 読 & -- \\
    監査履歴・固定層 & -- & -- & 追 & -- & -- & -- & -- & -- \\
    \bottomrule
  \end{tabularx}
  \caption{対象別の完全な疎権限行列。読、追、書、呼、有は、読出し、
    追記、書換え、呼出し、有効化を表す。空欄と退役指示文の読出しは
    すべて拒否する。実装は各役割・層・操作・対象を全展開した機械可読
    一覧も出力する。}
  \label{tab:authority-matrix}
\end{table}

候補の権限集合は、置換する現職の部分集合でなければならない。型付きの
実行時の型付き構築処理は、モデル本文中の作成者名や役割名を採用せず、現在の構成要素
識別子と固定役割を付与する。このため、モデル出力だけでは作成者を名乗れ
ない。ただし、これは同一処理内の実装上の境界であり、暗号署名、
役割ごとの鍵、OS 上の利用者分離、処理分離、処理間通信の隔離ではない。実行時、保存点の権限、
道具呼出し境界は現版の信頼基盤である。

\begin{table}[htbp]
  \centering
  \small
  \begin{tabularx}{\textwidth}{@{}p{0.28\textwidth}XX@{}}
    \toprule
    \textbf{攻撃者の能力} & \textbf{現実装が強制すること} &
    \textbf{保証しないこと} \\
    \midrule
    統治対象モデルの悪意ある文章 &
      型付き作成者付与、呼出し許可表、適法な遷移 &
      モデルが示す証拠の意味的な真実性 \\
    モデル役の結託または相関故障 &
      記録の分離と自己裁定の拒否 &
      統計的独立性や結託への耐性 \\
    同一処理内の役割コードの不具合・侵害 &
      通常の呼出し口を通る限り各検査を実行 &
      記憶領域やファイルへの直接接触の防止 \\
    保存点書き手またはホスト管理者 &
      敵対的保証なし &
      台帳、履歴、局所固定値、実行形式の完全性 \\
    同じモデル名の背後にある提供者側変更 &
      実行識別の不完全性を記録 &
      不透明なモデル重みの同一性 \\
    外部への不可逆操作 &
      本成果物には固定された外部操作門なし &
      装置、API、本番データに対する閉鎖側故障 \\
    \bottomrule
  \end{tabularx}
  \caption{脅威能力と保証の対応。現在の主張は、信頼されたホスト上の
  応用処理レベルの参照監視と作業手順完全性であり、OS や暗号による
  セキュリティ境界ではない。}
  \label{tab:threat-capability}
\end{table}

\subsection{出来事形式と再開手続き}\label{sec:appendix:events}

新しい出来事は形式版二を用い、単調増加する識別子、種類、取引識別子、
内容、直前ハッシュ値、時刻を持つ。完全長の SHA-256 値は
\[
  (\textit{形式版},\textit{出来事識別子},\textit{種類},
  \textit{取引識別子},\textit{内容},\textit{直前ハッシュ値})
\]
を正規化して結ぶ。したがって、種類、識別子、取引境界、直前参照、内容、
順序の変更で値が変わる。時刻は対象外の付随情報である。内容だけを結ぶ
旧版一も読出せるが、新規書込みは版二に限る。この値は作成者の署名では
ない。

\floatname{algorithm}{手続き}
\algrenewcommand\algorithmicrequire{\textbf{入力：}}
\algrenewcommand\algorithmicensure{\textbf{出力：}}
\algrenewcommand\algorithmicreturn{\textbf{返す}}
\begin{algorithm}[htbp]
\caption{再開時の決定論的な状態再構築}
\begin{algorithmic}[1]
\Require 保存順の出来事列 $E$、残存していれば直前境界の局所固定値 $P$
\Ensure 再構築した期状態と台帳、または完全性違反
\State 識別子が狭義単調増加することを検査する
\State 各出来事について版に応じたハッシュ値と直前参照を検査する
\State $P$ がある場合は、$E$ が $P$ の記録数以上であり、その位置の
ハッシュ値が $P$ の先頭値と一致することを検査する
\State 準備、構成出来事、確定の取引識別子が一致することを検査する
\State 準備の後に対応する確定がない末尾取引を除外する
\State 空の状態から確定済み出来事を順に畳み込み、期と台帳を再構築する
\State \Return 再構築状態
\end{algorithmic}
\end{algorithm}

版二は、再生の意味に関わる項目の改変、並べ替え、取引境界の差替え、
後続を残した途中削除、残存する局所固定値より短い履歴を検出する。書き手は
壊れた連鎖への追記を拒否し、再生は未確定の最終取引だけを事故回復として
除外する。しかし、履歴と局所固定値を同じ過去状態へ戻す末尾
切り詰めは検出できない。外部署名済みの最新先頭値、単調カウンタ、独立
監査者などは現実装にない。したがって、本稿の保証は単一保存点内部の
完全性に限り、管理者による保存点全体の巻戻しは脅威モデル外である。

\subsection{責任帰属、影響修復、原稿検証}\label{sec:appendix:algorithms}

\paragraph{責任帰属。}
生の攻撃は成果物だけを対象とする。裁定後、攻撃種別に対応する役割を
固定表から取り、その期に固定された現職へ解決する。各検証済み攻撃記録
は責任対象を一つだけ持つ。複数役割が関与する場合は、同じ裁定を参照する
役割別の記録を一件ずつ作る。一般研究向け七種の攻撃は、攻撃対象ノードに
作成時の構成要素、指示文、期が一度だけ付与され、それが期に固定された
生成役と一致する場合だけ生成役へ結び付ける。古い記録、欠落、不一致は
推測せず対象なしとする。論文の過剰受理は査読役へ結び、原稿が固定忠実性
検査にも失敗した場合は執筆役にも別記録を作る。複数原因は同じ裁定を参照
する複数の一対象記録で表す。この固定写像の因果的な正確性は未評価であり、
一対象形式は監査上の制限であって因果保証ではない。道具、データ、共有
指示文にまたがる原因はまだ表現していない。

\paragraph{影響修復。}
直接影響は、退役した指示文のハッシュ値を作成来歴に持つ記録である。
間接影響は、記録間の参照辺を逆向きに幅優先探索して求める。提案が過去
提案を着想資料として参照する一種類だけを非本質的な辺とし、未定義の辺は
保守的に本質的と扱う。到達閉包は既定深さ八を越えても最後まで求め、八を
越えた非本質辺の免除を止めて保守的に無効化する。到達記録数を $V$、逆参照
数を $E$ とすると、時間・記憶量は $O(V+E)$ であり、訪問済み集合で循環を
停止する。提案・振分け・効用規則の直接影響
はノードを選択対象外にし、査読・敵対検査・弁護・裁定の影響は残存入力
から点数を再計算する。必要な入力がなければ無効化する。この規則の意味的
な精度と過剰除去率は未評価である。

\paragraph{最終原稿検証。}
検証器は一般的な自然言語理解を行わない。原稿側が、主張識別子、数値
識別子、式、被演算値のノードと測定位置、単位、許容誤差を前向きに宣言
する。検証器は実行済みノードと成果物の存在を確かめ、閉じた式集合
（恒等、比、速度向上率、増減率、差）で数値を再計算し、原稿中の対応値
と絶対または相対誤差で比較する。要旨、結果、結論を厳格に走査し、序論、
考察、限界は警告対象とする。関連研究、参考文献、付録、数式は除外する。
再現可能な数値主張率と数値被覆率に加え、識別子衝突、未知の式、欠けた
被演算値、単位・実行環境の不一致を報告する。本稿は ARI の実行保存点
から生成したものではないため、実行証拠に根ざす有効な主張一覧は存在せず、
自己適用を主張しない。手作りの保存点で代用しても、提案する公開経路の
試験にはならない。
構造化宣言のない非数値的主張の真偽、実行記録自体の真実性、意味的な
同値表現は検証対象外である。したがって「固定」とは、この構造化入力に
対する処理がモデルを呼ばず決定的であることを意味し、自然言語全体を
正しく理解することを意味しない。

\subsection{決定性と検査証拠の境界}\label{sec:appendix:determinism}

本稿が保証する決定性は二つに限る。\emph{カーネル決定性}は、保存済みの
助言報告、候補評価、台帳、設定が同一なら、遷移と修復結果が同一になる
ことを指す。\emph{履歴再生決定性}は、確定済み出来事列が同一なら、派生
状態が同一になることを指す。言語モデルの弁護、裁定案、査読を同じ入力
から再生成できるとは主張しない。モデル提供者の版、重み、復号設定、
道具出力まで固定されていないためである。

本改訂の全体試験は 4,805 件を収集し、4,787 件が合格、16 件が条件付きで
省略、2 件が既知の失敗として扱われた。RQGM と最終原稿検証に直接関係
する 1,124 件も全件合格した。出来事の意味項目、直前参照、順序、
取引不一致、緊急境界の指紋値、権限一覧の完全性に対する負例を追加した。
重複しない内訳は、最終原稿検証 64、カーネル等 285、敵対検査・統治
126、修復・隔離再生成 66、効用進化 85、評価 157、指示文進化 93、
論文役 94、その他 154 件である。試験数は科学的な性能指標ではない。コード
被覆率と変異試験得点は測定しておらず、複数モデル、
複数乱数種による検出率・費用・研究品質の比較結果もまだない。
